%% kapa: exam=1 round=36 year=2025 period=1 form=A subject=real_estate subject_ko=부동산학원론 seq=12 number=92
%% source: Q-Net 감정평가사 기출문제 (cst00309) artlSeq=5231525 file=raw/1008/5231525/2025년 제36회 감정평가사 1차시험 1교시 문제지 원본.pdf
%% conversion: kordoc 4.24.51 -> tools/gen1.py (auto); review=proofread-sonnet; flags=none
\begin{question}{92}
어느 지역에 A점포와 B점포가 있다. A점포의 면적은 1,200 $\mathrm{m}^{2}$ 이고, B점포의 면적은 10,800 $\mathrm{m}^{2}$ 이다. A점포와 B점포 사이의 직선거리는 4 km이다. 컨버스(P. Converse)의 분기점 모형에 기초할 때, A점포와 B점포의 상권 경계지점은 B점포로부터 얼마만큼 떨어진 지점인가? (단, A점포와 B점포는 동일 직선상에 위치하며, 주어진 조건에 한함)
\begin{choices}
  \item 1 km
  \item 2 km
  \item 3 km
  \item 4 km
  \item 5 km
\end{choices}
\end{question}
\answer{3}
